\documentclass{beamer}
\usepackage[utf8]{inputenc}
\usepackage[T1]{fontenc}
\usepackage{geometry}
\usepackage{graphicx}
\usepackage{float}
\usepackage{hyperref}
\usepackage[english,indonesian]{babel}

\title{Optimasi Pemrosesan Template LaTeX Menggunakan Pendekatan Berbasis AI}
\author{Academic Support Agent}
\date{\today}

\begin{document}
\maketitle

\begin{abstract}

\end{abstract}

\chapter{PENDAHULUAN}


\chapter{TINJAUAN PUSTAKA}


% Original content (commented out):
% %---------------------------------------------------
% Template para slides com as cores da UFV
% Autor.: Rodrigo Smarzaro (smarzaro@ufv.br)
% Versão: 1.0 
%---------------------------------------------------
% As configurações estão no documento preambulo.tex
% Coloque todas as figuras na pasta \Figuras
% Deve ser compilado com pdfLaTeX + Biber
%---------------------------------------------------

\documentclass{beamer}
% use a versão abaixo para desativar as transições.
%\documentclass[handout]{beamer}

%carrega os ajustes padrões para os slides
\input{preambulo}


\title{Template para Slides - UFV}
\subtitle{Curso de Sistemas de Informação}
% Você pode alterar o logo usado no arquivo preambulo.tex
\author[Seu nome aqui]{\texorpdfstring{\logo \\ Seu Nome\\ \texttt{Seu e-mail}}{Seu Nome}}
\institute[UFV-CRP]{Universidade Federal de Viçosa}
\date{2021}

\begin{document}

% Para controlar se os tópicos aparecem com transparência ou invisível
\transparente
%\invisivel

%Slide de Título
\frame{\titlepage}


\section*{Conteúdo}

\section{Apresentação}

%exemplo de um slide
\begin{frame}{Exemplos de blocos}
	%\framesubtitle{subtitle}
\begin{block}<+->{Exemplo de um bloco normal (\alert{\texttt{block}})}
	Blocos são úteis para destacar uma parte do conteúdo no slide.
\end{block}
\vfill
\begin{exampleblock}<+->{Exemplo de um bloco ``exemplo'' (\texttt{exampleblock})}
	Blocos são úteis para destacar uma parte do conteúdo no slide. 
\end{exampleblock}
\vfill	
\begin{alertblock}<+->{Exemplo de um bloco ``alerta''(\alert{\texttt{alertblock}})}
	Blocos são úteis para destacar uma parte do conteúdo no slide. 
\end{alertblock}
\vfill
Cores\footnote{Veja aqui as opções de configuração de visual do beamer \url{https://tinyurl.com/beamerconfigs}} podem ser alteradas no arquivo \texttt{preambulo.tex}
\end{frame}

\begin{frame}{Listas}
	\begin{itemize}[<+->]
		\item Listas
		\vfill
		\item Com vários itens\dots
		\begin{itemize}
			\item e subitens
		\end{itemize}
	\end{itemize}
\end{frame}

\begin{frame}[fragile]{Código-fonte}
	\begin{itemize}
		\item Você pode incluir facilmente código-fonte nos slides com o pacote \texttt{listings}
		\begin{lstlisting}
CREATE TABLE DEPARTAMENTO(	
DNumero    INTEGER     NOT NULL,
DNome      VARCHAR(10) NOT NULL,
CPFGerente CHAR(11),
Datainicio DATE)
		\end{lstlisting}    

		\item Se fizer isto (e caso uso o ambiente \texttt{verbatim} o slide deve ser declarado como:

    \begin{lstlisting}[style=TeX-Smarzaro]
\begin{frame}[fragile]{Nome do Slide}
    \end{lstlisting}		
	\end{itemize}
\end{frame}

\begin{frame}{Referências}
    % incluir as principais referências aqui caso não tenham sido citadas explicitamente
    % nos slides
    \nocite{Silberschatz2014,Elmasri2011,Coronel2010,Date2004,Teorey2013,Puga2013}
    \printbibliography
\end{frame}

\end{document}


\end{document}